

Upon introducing the multiparameter stochastic sewing lemma and the notions of additive and multiplicative local non-determinism, we were able to establish new regularity estimates for the local times of Gaussian fields. We observed that additive LND yields local times whose regularity benefits from each parameter in an additive fashion, while the regularity of local times associated to multiplicatively LND Gaussian fields is essentially dictated by the "weakest" parameter. Moreover, in Remark \ref{boundary remark} we saw that additive LND can only be obtained from the boundary of a stochastic field, meaning in particular that processes with deterministic boundary (such as the fractional Brownian sheet) can not be additively LND.  Unfortunately, we are not able to show that our regularity results for the local times are sharp. However, Remark \ref{boundary remark} is rather suggestive of the following conjecture: 
\begin{con}
    For a Gaussian stochastic field, the maximal regularity is in general dictated by its behaviour on the boundary. 
\end{con}

\begin{defn}[Strong LND] \cite[Section 2.4, (C3')]{Xiao} Let $\zeta\in \RR^d_+$. A Gaussian random field $\{W_\bt\}_{\bt\in \RR^n}$ with stationary increments admits the $\zeta$-strong local non-determinism property if there exists a constant $c>0$ such that for any $m\geq 1$ and $\bu, \bt^1, \dots \bt^m\in \RR^d$ we have
\begin{equation*}
    \mbox{Var}(W_\bu\ |\ W_{\bt^1}, \dots W_{\bt^m})\geq  c\min_{1\leq k\leq m}\left(\sum_{i=1}^d |\bt^k_i-\bu_i|^{\zeta_i}\right)^2. 
\end{equation*}
\end{defn}
 \begin{cor}
        A continuous Gaussian process with commuting filtrations that admits the $\zeta$-strong local non-determinism property satisfies for $\bt>\bs$
        \begin{equation}
            \mbox{Var}(W_\bt| \mathcal{F}_\bs)\geq c\sum_{i=1}^d |t_i-s_i|^{2\zeta_i}
            \label{strong LND impli}
        \end{equation}
    \end{cor}
    \begin{proof}
        This follows from a simple adaption of the  proof in Corollary \ref{cor:sectorial_non_determinism}. 
    \end{proof}
An example that admits the strong local non-determinism property is given by Gaussian processes whose spectral density decay can be controlled in a quantified way.
\begin{defn}[Spectral density of Gaussian processes]
    Let $(X_\bt)_{\bt\in \RR^d}$ be a centered Gaussian processes starting in zero with stationary incerements and continuous covariance function $R(\bs, \bt)=\mathbb{E}[X_\bs X_\bt]$. Then, by Yaglom \cite{Yaglom1957} there exists a non-negative definite matrix $\Theta$ and a non-negative symmetric measure $\mu(d\xi)$ on $\RR^d\backslash\{0\}$ satisfying $\int_{\RR^d}\frac{|\xi|^2}{1+|\xi|^2}\mu(d\xi)<\infty$ such that 
    \[
    R(\bs, \bt)=\int_{\RR^d}(\exp(i\langle \bs, \xi\rangle)-1)(\exp(-i\langle \bt, \xi\rangle)-1)\mu(d\xi)+\langle \bs, \Theta \bt\rangle
    \]
    The measure $\mu$ is called spectral measure of $X$. If $\mu$ admits a density $f$ with respect to the Lebesgue measure on $\RR^d$, $f$ is called spectral density of $X$. 
\end{defn}

In the next theorem we recall a sufficient condition on the spectral density associated to a Gaussian field which guarantees that the field is $\zeta$-strong LND. 

\begin{thm}\cite[Theorem 3.2]{Xiao}
    Let $(X_\bt)_{\bt\in \RR^d}$ be a centered Gaussian processes starting in zero with spectral density $f$. Assume there exists $\zeta\in (0,1)^d$ such that 
    \[
    f(\xi)\simeq \frac{1}{\left(\sum_{i=1}^d|\lambda_i|^{\zeta_i}\right)^{2+Q}}, \qquad \lambda\in \RR^d\backslash \{0\}
    \]
    where $Q=\sum_{i=1}^d\frac{1}{\zeta_i}$. Then $X$ admits the $\zeta$-strong local non-determinism property.
\end{thm}

\begin{rem}
    If spectral measure factorizable, then we might preserve commuting filtrations (Hannes?) Hence, we could use Xiaos result on strong local non-determinism again. However, these seem mutually exclusive: If spectral measure is factorizable, it's growth will not be similar to the one demanded in Theorem 3.2 of Xiao, which imposes an additive structure. 
\end{rem}




\begin{rem}
i.e. we shift our fractional Brownian sheet away from zero, meaning we can use Theorem 3.1 in (Xiao, SPDE lecture notes), i.e. we have (sectorial) local non-determinism, which is enough for our purposes. (However: slnd obtained through Fourier arguments, so how do we get commuting filtrations again?)  Question: Why is the boundary pathological? For Riemann-Liouville that is clear, but what about fractional Brownian sheet?
\end{rem}



\subsection{The Riemann-Liouville fractional Brownian sheet}

We will now define a type of fractional Brownian sheet, inspired by the Riemann Liouville fractional Brownian motion. This random field will work as the main example throughout the article. As the Brownian sheet in itself is zero on the boundary, we include boundary terms in the definition, and we will in later applications see the importance of this terms if chosen appropriately.

\begin{defn}
Let $H=(H_1,\ldots,H_N)\in (0,1)^d$ and consider a Brownian sheet $W:\Omega\times [0,T]^d\rightarrow \RR$ on a probability space $(\Omega,(\cF_{\bt})_{\bt},\PP)$. For any proper subset  $\eta\subset [d]$,   $\xi^\eta:\Omega\times [0,T]^{|\eta|} \rightarrow \RR$ represents stochastic boundary terms adapted to the filtration $(\cF_{\bt})_{\bt}$. 
We define the multiparameter Riemann Liouville (RL) fractional Brownian sheet with the stochastic boundary terms $\xi^\eta$ as the following multiparameter Wiener-Walsh integral for any $\bt\in [0,T]^d$
\begin{equation}
    W^H_\bt = \sum_{\eta \subsetneq [d]} \xi_{\bt_{\eta}}^\eta + \int_\boo^\bt K(\bt,\bs) \dd W_\bs.
\end{equation}
where $K(\bt,\bs)= \prod_{i=1}^d |t_i-s_i|^{H_i-\frac{1}{2}}$, and $\bt_\eta=(t_{\eta_1},...,t_{\eta_{|\eta|}})$.
\end{defn}

\begin{rem}
Note that if $\xi^\eta=0$ for all $\eta\subset \theta$ this implies that the fractional Brownian sheet is $0$ on the boundary of the rectangle $[0,T]^2$.  We include the boundary terms $\xi^i$ for $i=1,2$ in the definition to provide sufficient flexibility towards applications.  
\end{rem}





In the next results we will limit the analysis to the case $d=2$ for notational ease and brevity. We believe this is more instructive, and the techniques are readily extendable to the higher dimensional case. 
 
\begin{prop}\label{prop:conditional variance}
 Let $W^H:[0,T]^2\times \Omega \rightarrow \RR$ be a RL fractional Brownian sheet with $H\in (0,1)^2$  on a probability space $(\Omega, \cF,\cF_{\bt},\PP)$ where $\cF_\bt$ is a commuting filtration. The boundary terms is given by $\xi^1,\xi^2:[0,T]\times \Omega \rightarrow \RR$, two real valued independent stochastic processes adapted to the filtration $\{\cF_\bt\}$. Then for any $\bs\leq \bt$ we have 

\begin{multline}\label{eq:conditional variance}
    Var(W_\bt|\cF_\bs) = 
    Var(\xi^1_{t_1}|\cF_\bs)+Var(\xi^2_{t_2}|\cF_\bs)
    \\
    +
    c_H[|t_1-s_1|^{2H_1}(s_2^{2H_2}+\frac{1}{2}|t_2-s_2|^{2H_2})+|t_2-s_2|^{2H_2}(s_1^{2H_1}+\frac{1}{2}|t_1-s_1|^{2H_1})]   
\end{multline}
for some constant $c_H>0$. 
\end{prop}

\begin{proof}
We write the proof for $N=2$ but the same proof can be adapted to general $N$. First, recall that 
\begin{equation*}
    Var(W^H_\bt |\cF_\bs)=\EE[(W^H_\bt -\EE[W^H_\bt |\cF_\bs])^2|\cF_\bs],  
\end{equation*}
and observe that by elementary algebraic manipulations we have  
\begin{equation*}
\begin{aligned}
    W^H_{t} &= \int_{0}^{\bs}  K(\bt,\bu)\dd W_\bu  + \int_{0}^{s_1}\int_{s_2}^{t_2} K(\bt,\bu)\dd W_\bu + \int_{s_1}^{t_1}\int_{0}^{s_2} K(\bt,\bu)\dd W_\bu +\int_{s_1}^{t_1}\int_{s_2}^{t_2}  K(\bt,\bu)\dd W_\bu 
    \\
    &=:\int_{0}^{\bs}  K(\bt,\bu)\dd W_\bu +I_1 +I_2+I_3. 
    \end{aligned}
\end{equation*}

Furthermore,  note that since $W$ on $[s_1,t_1]\times [0,s_2]$ is independent of $\cF_{s_1,s_2}$ (See illustration Figure \ref{filtrations}), the integral $I_1$ is independent from $\cF_{s_1,s_2}$. Similarly, we argue that also $I_2$ and $I_3$ are independent of $\cF_{s_1,s_2}$. 
In addition, it is readily checked that 
\[
\EE[W^H_\bt |\cF_\bs]= \EE[\xi^1_{t_1}|\cF_\bs]+\EE[\xi^2_{t_2}|\cF_\bs]+\int_{0}^{\bs}  K(\bt,\bu)\dd W_\bu 
\]
Inserting this into the expression for the conditional variance, using independence between $\xi^1,\xi^2 $ and $W$, we see that we only need to compute $\EE[I_iI_j|\cF_\bs]$ for all $i,j=1,2,3$.  Suppose now that  $i\neq j$, and then we see that 
\begin{equation*}
    \EE[I_iI_j|\cF_\bs]=0 \
\end{equation*}
and for $i=1,2,3$ and $j=1,2$ we have, due to independence, 
\[
 \EE[I_i (\xi^j_{t_j}-\EE[\xi^j_{t_j}|\cF_\bs])]=0. 
\]
%
Therefore, since $I_1,I_2,I_3$ are all independent of $\cF_\bs$ we have that 
\begin{equation*}
    \EE[(W_\bt -\EE[W_\bt |\cF_s])^2|\cF_\bs]=\sum_{i=1}^3\EE[I^2_i]. 
\end{equation*}
By the It\^o isometry, we then get that 
\begin{equation*}
    \EE[I^2_1] =\int_{0}^{s_1} \int_{s_2}^{t_2} K(t,u)^2 du_1du_2 
\end{equation*}
Using the explicit representation of the Riemann-Liouville kernel, we compute  
\begin{equation*}
    \EE[I^2_1]=\frac{s^{2H_1}_1}{2H_1}\frac{(t_2-s_2)^{2H_2}}{2H_2}, \quad \EE[I^2_2]=\frac{s^{2H_2}_2}{2H_2}\frac{(t_1-s_1)^{2H_1}}{2H_1} \quad \EE[I^2_3]=\frac{(t_2-s_2)^{2H_2}}{2H_2}\frac{(t_1-s_1)^{2H_1}}{2H_1}. 
\end{equation*}
And thus, summing these terms we get 
\begin{equation*}
    \sum_{i=1}^3 \EE[I^2_i]= \frac{1}{4H_1H_2}\left((t_1-s_1)^{2H_1}(s_2^{2H}+\frac{1}{2}(t_2-s_2)^{2H_2})+(t_2-s_2)^{2H_2}(s_1^{2H_1}+(t_1-s_1)^{2H_1})\right). 
\end{equation*}
Combining the above computation provides the conclusion in \eqref{eq:conditional variance} with $c_H=(4H_1H_2)^{-1}$. 


\end{proof}

In the following corollary, we compute the conditional variance of the fractional Brownian sheet motion with a boundary given by two independent Riemann Liouville fBm's.  



\begin{cor}
Let $\xi^1,\xi^2:[0,T]\times \Omega \rightarrow \RR$ be two real valued independent RL-fractional Brownian motions with Hurst parameter $H_1,H_2\in(0,1)$. Let $W^H$ be a RL fractional Brownian sheet with $H\in (0,1)^N$ and boundary terms $\xi^i$ for $i=1,2,$ on a probability space $(\Omega, \cF,\cF_{\bt},\PP)$ where $\cF_\bt$ is a commuting filtration. Then for any $\bs\leq \bt$ we have 

\begin{equation*}
    Var(W_\bt|\cF_\bs) = c_H[|t_1-s_1|^{2H_1}(s_2^{2H_2}+\frac{1}{2}|t_2-s_2|^{2H_2}+ 2H_2)+|t_2-s_2|^{2H_2}(s_1^{2H_1}+\frac{1}{2}|t_1-s_1|^{2H_1}+ 2H_1)].
\end{equation*}
In particular, there exists a constant $c>0$ such that 
\[
Var(W_\bt|\cF_\bs) \geq c\sum_{i=1}^d |t_i-s_i|^{2H_i}
\]
\end{cor}

\begin{proof}
    Since $\xi^i$ is an RL fBm, we have for $i=1,2,$
    \begin{equation}
        Var(\xi^2_{t_2}|\cF_\bs)=Var(\xi^i_{t_i}|\cF^i_{t_i})=\frac{1}{2H_i}(t_i-s_i)^{2H_i}. 
    \end{equation}
    Inserting this into the relation obtained in Proposition \ref{prop:conditional variance} prove the claim. 
\end{proof}